\section{Chapter 14}

\begin{verse}
One looks at it and does not see it, so that it is called invisible\\
One hears it and does not understand it, so that it is called soundless\\
One touches it and does not grab it, so it is called incorporeal\\
These three characteristics of the transcendent\\
Mix together and say one thing.\\
It is the negation of the higher luminosity\\
And of the lower darkness\\
It is immense and insusceptible of name\\
It returns to what stands before (and above) being\\
Its form is the absence of form\\
Its figure is the absence of figure\\
Undiscernible profundity\\
If you look ahead, you will not see the end\\
If you look behind, you will not see the origin\\
Primordial principle\\
In action now as in every time = perennially current \\
One follows the way of the Ancients\\
And one will know the eternal essence of the Principle.
\end{verse}

The ninth line can be understood both metaphysically: The Tao is before and above being --- as well as in these terms, referring to individual beings: the Tao brings them back from the current of transformations to that which is before and above being (in the pre-formal state).

In the first lines: the Principle, tangible but elusive. Non-being, which is not the not-existent. This passage of Kuan Tze (XIII, 36) will be interesting, for the meeting of the attributes of the principle with the technique of realization:

\begin{quotex}
The Tao is not far, yet reaching it is difficult ... Empty and non-being, incorporeity and immateriality, that is called Tao. Heaven is empty, Earth is calm, neither one nor the other struggles. The I is discarded and falls silent: then divine clarity will persist in you. Who understands silence (\emph{p'u jen}) and immobility profoundly will understand the essential drama of the Tao... If (man) will renounce craving, the emptiness (\emph{hsu}) will entirely compenetrate him; compenetrated by it, he will remain calm and non-acting; calm and firm, he will take up contact with the vital ether; whoever possesses the vital ether is detached (from terrestrial form); detached, he will radiate light like a god (\emph{ming shen}).
\end{quotex}


\textbf{Chinese text and literal translation}


\subsubsection{Chapter 14 (\begin{CJK*}{UTF8}{bsmi}第十四章\end{CJK*})}
\columnratio{.4}
\begin{paracol}{2}
\begin{verse}
\begin{CJK*}{UTF8}{bsmi}
視之不見，名曰夷；\\
聽之不聞，名曰希；\\
摶之不得，名曰微。\\
此三者不可致詰，故混而為一。\\
其上不皦，其下不昧，\\
繩繩不可名，復歸於無物，\\
是謂無狀之狀，無物之象，\\
是謂惚恍。\\
迎之不見其首，隨之不見其後。\\
執今之道以御今之有，能知古始，\\
是謂道紀。
\end{CJK*}
\end{verse}
\switchcolumn
\begin{verse}
That which is looked at yet never seen is known as invisible,\\
That which is listened to yet never heard is known as soundless,\\
And that which is grasped yet never touched is known as intangible.\\
The origin of these three can not be observed, so mix them and form a unity.\\
Above, it is not lit; Below, it is not dark.\\
The Dao cannot be named by common rules, as it returns to its origin,\\
It is the shape of the original shape, the image of the original object,\\
This is blur and profound.\\
Step forward, cannot see its beginning; Follow it, cannot see its end.\\
Hold the ancient Dao to know the present things, and thus know the ancient beginnings,\\
This is the principle of the Dao.
\end{verse}
\end{paracol}

\flrightit{Posted on 2020-11-17 by Lao Tzu}

